% sup1_04_twodomain: Appendix on Ba3 Prop. 1, part 1, step S2: delta G_k(Omega, Omega_2, B~) as an external field (first alternative of (2.5)).

\subsubsection{S2: the two-domain propagator $\delta G_k(\Omega,\Omega_2,\tilde B)$ as an external field}\label{sup1:ssec:S2}

\cite[p.~414]{Ba3} replaces $\Omega$ by $\Omega_2\subset\Omega$ ``satisfying the same conditions as $\Omega$'', and treats $\delta G_k(\Omega,\Omega_2,\tilde B)=G_k(\Omega,\tilde B)-G_k(\Omega_2,\tilde B)$ as an external field; \cite[(2.5)]{Ba3}, first alternative, asserts $\norm{h\,\delta G_k(\Omega,\Omega_2,\tilde B)\,h'}_{1,\alpha}\le O(e^{-\delta_0\dist(\Omega_2,\partial\Omega)})e^{-\delta_0\dist(\supp h,\supp h')}$, and \cite[p.~426]{Ba3} uses it in \cite[(2.13)]{Ba3} as an $O(1)$ bound with the two-point decay; the small depth factor is what the domain changes of \cite[pp.~421--422]{Ba3} use. $\delta G$ is defined on $\Omega_2\times\Omega_2$, $G_k(\Omega_2,\tilde B)$ being extended by $0$.

\paragraph{The reading of the depth factor.} The printed factor $e^{-\delta_0\dist(\Omega_2,\partial\Omega)}$, read as a distance between the sets $\Omega_2$ and $\partial\Omega$, cannot be what is meant: $\delta G$ is of order one at points near $\partial\Omega_2$ whatever the distance to $\partial\Omega$. (For $\Omega=T_\eta$ there is no $\partial\Omega$ at all; in one dimension at $\eta=\frac15$, $A=0$, $m=0$, $a_k=1$, with $\Omega_2$ a union of $20$ blocks, $\delta G(x,x)$ takes the values $-0.43$, $-0.28$, $-0.05$, $-0.0008$ at depth $0$, $0.2$, $1$, $3$ in $\Omega_2$, against $G(x,x)=0.52$; the same values arise with $\Omega$ a Neumann segment at distance $90$ from $\Omega_2$.) The factor holds in the form of \cite[Prop.~2.1, (2.26)]{Ba1}, where the extra decay is in the depth of the arguments in the \emph{smaller} set: here, the depth of the vertex supports in $\Omega_2$. This is the form proved below, and it is the form the uses need.

\paragraph{Setting.} $X_{\mathrm{in}}\subset X_{\mathrm{out}}$ are unions of large blocks of $T_\eta$, $X_{\mathrm{out}}$ may be $T_\eta$, $X_{\mathrm{in}}\ne T_\eta$; $A$ is real and regular by \eqref{sup1:eq:E}, $0\le m^2t^2\le c_m$; $H_{\mathrm{in/out}}\defeq H_{X_{\mathrm{in/out}}}(A)$, $G_{\mathrm{in/out}}\defeq H_{\mathrm{in/out}}^{-1}$, $\delta\defeq G_{\mathrm{out}}-G_{\mathrm{in}}$, $\rho\defeq\dist_\infty(\cdot,X_{\mathrm{in}}^c)$; then $\rho_{X_{\mathrm{out}}}\ge\rho$.

\begin{lemma}[collars]\label{sup1:lem:S21}
Fix $f\in C^\infty(\R)$ with $0\le f\le1$, $f=0$ on $(-\infty,0]$, $f=1$ on $[1,\infty)$, and $\varphi\in C^\infty_c([-\frac14,\frac14]^3)$, $\varphi\ge0$, $\int\varphi=1$. For a unit block $b$ with lowest corner $y_b$ put $\hat b\defeq y_b+[0,1)^3$ and $\hat X_{\mathrm{in}}\defeq\bigcup_{b\subset X_{\mathrm{in}}}\hat b$; let $\rho^c$ be the continuum sup-distance to the complement of $\hat X_{\mathrm{in}}$, $\tilde\rho\defeq\rho^c*\varphi$, and $\chi_s\defeq f(\tilde\rho-s)$ at lattice points; $c_1'\defeq\norm{f'}_\infty$, $c_2'\defeq\norm{f''}_\infty+\norm{f'}_\infty\max_\nu\norm{\partial_\nu\varphi}_{L^1}$. Then: \textup{(i)} $\tilde\rho$ is $1$-Lipschitz for $|\cdot|_\infty$ and $\rho-\frac14-\eta\le\tilde\rho\le\rho+\frac14$ at lattice points; \textup{(ii)} $\chi_s=0$ on $\{\rho\le s-\frac14\}$ and $\chi_s=1$ on $\{\rho\ge s+\frac54+\eta\}$; \textup{(iii)} $|d^\eta_\mu\chi_s|\le c_1'$ and $|d^\eta_\mu d^\eta_\nu\chi_s|\le c_2'$, uniformly in $\eta$, $s$, $X_{\mathrm{in}}$; \textup{(iv)} if $\chi_s(u)\ne\chi_s(u')$ and $|u-u'|_\infty\le l$, then $\tilde\rho(u),\tilde\rho(u')\in(s-l,s+1+l)$.
\end{lemma}

\begin{proof}
(i) A lattice point outside $X_{\mathrm{in}}$ lies in a cube $\hat b'$ with $b'\not\subset X_{\mathrm{in}}$, so $\rho^c\le\rho$; the infimum defining $\rho^c(z)$ is attained on the closure of such a cube, whose lattice points are $\eta$-dense for $|\cdot|_\infty$, so $\rho\le\rho^c+\eta$; the mollifier preserves the Lipschitz constant and moves values by at most $\frac14$. (ii) follows from (i). (iii) Lattice differences are averages of derivatives along segments; $|\partial_\mu\tilde\rho|\le1$ and $\partial_\mu\partial_\nu\tilde\rho=(\partial_\mu\rho^c)*\partial_\nu\varphi$. (iv) If $\tilde\rho(u)\le s-l$ then $\tilde\rho(u')\le s$ and both values of $\chi_s$ are $0$; symmetrically both are $1$ if $\tilde\rho(u)\ge s+1+l$.
\end{proof}

\begin{lemma}[commutators]\label{sup1:lem:S22}
Let $s\ge1$, $\chi\defeq\chi_s$, $\phi_\mu\defeq d^\eta_\mu\chi$. \textup{(i)} $H_{\mathrm{out}}(\chi g)=H_{\mathrm{in}}(\chi g)$ for every $g$ on $X_{\mathrm{in}}$ (extended by $0$). \textup{(ii)} $C_\chi\defeq H_{\mathrm{out}}\chi-\chi H_{\mathrm{in}}=[H_{\mathrm{in}},\chi]=-\mathfrak r_\chi\mathfrak D=\mathfrak D^*\mathfrak r_\chi^*$, where $\mathfrak r_0\defeq\sum_\mu(d_\mu\phi_\mu)(\cdot-\eta e_\mu)-a_k[P_k(A),\chi]$ and $(\mathfrak r_\mu g)(x)\defeq\phi_\mu(x)g(x)+\phi_\mu(x-\eta e_\mu)U_{x-\eta e_\mu,\mu}(A)^{-1}g(x-\eta e_\mu)$. \textup{(iii)} $\sum_{i=0}^3\sup_z\sum_w\eta^d\norm{\mathfrak r_i(w,z)}\le c_{\mathfrak r}\defeq6c_1'+3c_2'+ac_1'$. \textup{(iv)} If $\mathfrak r_i(w,z)\ne0$ then $|w-z|_\infty\le1$ and $w,z\in S_s\defeq\{s-\frac54-\eta<\rho<s+\frac94+2\eta\}$.
\end{lemma}

\begin{proof}
(i) If $\chi(y)\ne0$ then $\rho(y)>s-\frac14\ge\frac34>\eta$, so every bond at $y$ and the unit block of $y$ lie in $X_{\mathrm{in}}$; a bond of $X_{\mathrm{out}}$ which is not a bond of $X_{\mathrm{in}}$ has both ends at depth $\le\eta$, outside $\supp\chi$; $P_k$ acts inside unit blocks. (ii) $D_\mu(\chi f)=\chi(\cdot+\eta e_\mu)D_\mu f+\phi_\mu f$ and $D_\mu^*(\chi(\cdot+\eta e_\mu)g)=\chi D_\mu^*g-\phi_\mu g$, so $[D_\mu^*D_\mu,\chi]=D_\mu^*(\phi_\mu\,\cdot)-\phi_\mu D_\mu$, and $D_\mu^*(\phi_\mu f)(x)=\eta^{-1}(\phi_\mu(x-\eta e_\mu)U^{-1}f(x-\eta e_\mu)-\phi_\mu(x)f(x))=-\phi_\mu(x-\eta e_\mu)U^{-1}_{x-\eta e_\mu,\mu}(D_\mu f)(x-\eta e_\mu)-(d_\mu\phi_\mu)(x-\eta e_\mu)f(x)$ (\cite[(2.32)]{Ba4}); collecting terms gives $-\mathfrak r_\chi\mathfrak D$, and $C_\chi^*=-C_\chi$ since $H_{\mathrm{in}}$ is self-adjoint and $\chi$ real. (iii) The two parts of each $\mathfrak r_\mu$ give $2c_1'$, so $6c_1'$ over three directions; the multiplication part of $\mathfrak r_0$ gives $3c_2'$; $[P_k,\chi](w,z)=P_k(w,z)(\chi(z)-\chi(w))$ with $|\chi(z)-\chi(w)|\le c_1'|\tilde\rho(z)-\tilde\rho(w)|\le c_1'$ and $\sum_w\eta^d\norm{P_k(w,z)}\le1$. (iv) \Cref{sup1:lem:S21}(iv) with $l=2\eta$ (bond and multiplication parts) or $l=1-\eta$ (block part) gives $\tilde\rho\in(s-1+\eta,s+2-\eta)$; convert to $\rho$ by \cref{sup1:lem:S21}(i).
\end{proof}

\paragraph{Choices.} $\chi\defeq\chi_{R_0+3}$, $\psi\defeq\chi_{R_0+13/2}$, $\mathcal P_D\defeq\{z:\rho(z)\ge R_0+10+2\eta\}\supset\{\rho\ge R_0+10.4\}$. With $\eta\le\frac15$, \cref{sup1:lem:S21,sup1:lem:S22} give: $S_\chi\subset\{R_0+1.55<\rho<R_0+5.65\}$, $S_\psi\subset\{R_0+5.05<\rho<R_0+9.15\}$; $\psi=0$ on $\{\rho\le R_0+6.25\}\supset S_\chi$; $\chi=1$ on $\{\rho\ge R_0+4.45\}$ and $\psi=1$ on $\{\rho\ge R_0+7.95\}$, both containing $\mathcal P_D$ and every bond partner of a point of $\mathcal P_D$; for $z\in\mathcal P_D$, $\dist_\infty(z,S_\chi)>\max\{\rho(z)-R_0-7,4.75\}$ and $\dist_\infty(z,S_\psi)>\max\{\rho(z)-R_0-10,1.25\}$; every point of $S_\chi$, $S_\psi$ and its bond partners has depth $>R_0+1.5$ in $X_{\mathrm{in}}$, hence in $X_{\mathrm{out}}$.

\begin{proposition}[the two-domain theorem]\label{sup1:prop:S23}
Let $0<\alpha<1$, $0<\beta\le\beta_D\defeq\min\{\beta_*,\delta_K/4\}$, $a,b\in\{0,1\}$, $Y\defeq D^a\delta D^{b*}$. For $x,x'\in\mathcal P_D$ with H\"older partners in $\mathcal P_D$,
\[
e^{\beta|x-x'|_\infty}\max\{\norm{Y(x,x')},aH_x,bH_{x'},abH_{xx'}\}\le C_D(\alpha)\exp\!\Big(-\frac{\delta_K}2\big[(\rho(x)-R_0-7)+(\rho(x')-R_0-10)\big]\Big),
\]
with $C_D(\alpha)\defeq c_Wc_{\mathrm{end}}(\alpha)^2$, $c_{\mathrm{end}}(\alpha)\defeq c_{\mathfrak r}e^{2\delta_K}\max\{c_K,2c_K(\alpha)\}c_3'^{1/2}$, $c_3'\defeq\sup_\eta\sum_z\eta^de^{-(\delta_K/2)|z|_\infty}\le(\frac15+4/\delta_K+1)^3$. The constants depend only on $d,a,M,L,\gamma_0,K_W,c_K,c_K(\alpha),f,\varphi$, not on $\eta,k,\eps$, the torus side, $X_{\mathrm{in}}$, $X_{\mathrm{out}}$, $A$, $m^2t^2\in[0,c_m]$ or $R_0$.
\end{proposition}

\begin{proof}
$G_{\mathrm{out}}C_\chi G_{\mathrm{in}}=\chi G_{\mathrm{in}}-G_{\mathrm{out}}\chi$, so at $y,y'$ with $\chi(y)=\chi(y')=1$, $\delta(y,y')=-(G_{\mathrm{out}}C_\chi G_{\mathrm{in}})(y,y')$, the inner variable ranging over $S_\chi$ only (\cref{sup1:lem:S22}(iv)), where $\psi=0$. Likewise $G_{\mathrm{in}}C_\psi G_{\mathrm{in}}=\psi G_{\mathrm{in}}-G_{\mathrm{in}}\psi$ gives $G_{\mathrm{in}}(z,y')=-(G_{\mathrm{in}}C_\psi G_{\mathrm{in}})(z,y')$ when $\psi(z)=0$, $\psi(y')=1$. Inserting and using \cref{sup1:lem:S22}(ii), at all points of $\mathcal P_D$ and their bond partners (the only points at which $D^a$, $D^{b*}$ evaluate $\delta$),
\[
\delta=-G_{\mathrm{out}}\mathfrak r_\chi(\mathfrak DG_{\mathrm{in}}\mathfrak D^*)\mathfrak r_\psi^*G_{\mathrm{in}},\qquad Y(x,x')=-\ip{E_L(x,\cdot)}{\mathfrak ME_R(\cdot,x')},
\]
$E_L\defeq D^aG_{\mathrm{out}}\mathfrak r_\chi$, $\mathfrak M\defeq\mathfrak DG_{\mathrm{in}}\mathfrak D^*$, $E_R\defeq\mathfrak r_\psi^*G_{\mathrm{in}}D^{b*}$. \emph{Left end.} $E_L(x,z)=\sum_w\eta^dD^a_xG_{\mathrm{out}}(x,w)\mathfrak r_\chi(w,z)$ with $w,z\in S_\chi$, $|w-z|_\infty\le1$; (K) applies on $X_{\mathrm{out}}$ at $(x,w)$ (both deep, $|x-w|_\infty\ge4.75\ge1$), so $K^{\delta_K}_{2-a}(x-w)\le e^{-\delta_K|x-w|_\infty}$ and $\norm{E_L(x,z)}\le c_{\mathfrak r}c_Ke^{\delta_K}e^{-\delta_K|x-z|_\infty}\mathbf 1_{S_\chi}(z)$; with $2(\delta_K-\beta)\ge\frac32\delta_K$, $\norm{e^{\beta|\cdot-x|}E_L(x,\cdot)}_2\le c_{\mathfrak r}c_Ke^{\delta_K}e^{-(\delta_K/2)\dist(x,S_\chi)}c_3'^{1/2}\le c_{\mathrm{end}}e^{-(\delta_K/2)(\rho(x)-R_0-7)}$; for $a=1$ and a partner $\tilde x\in\mathcal P_D$, (K2) with $|\tilde x-w|\ge1$ gives the transported difference with the factor $2c_K(\alpha)e^{\delta_K}|x-\tilde x|^\alpha$. \emph{Right end.} $(G_{\mathrm{in}}D^{b*})(v,x')=(D^b_{x'}G_{\mathrm{in}}(x',v))^T$ with $v\in S_\psi$, and (K) on $X_{\mathrm{in}}$ at $(x',v)$ gives $\norm{e^{\beta|\cdot-x'|}E_R(\cdot,x')}_2\le c_{\mathrm{end}}e^{-(\delta_K/2)(\rho(x')-R_0-10)}$, and the same times $|x'-\tilde x'|^\alpha$ for the transported difference. \emph{Middle.} \Cref{sup1:lem:W2} on $X_{\mathrm{in}}$ with weight $|\cdot-x|_\infty$ gives $\norm{e^{-\beta|\cdot-x|}\mathfrak Me^{\beta|\cdot-x|}}\le c_W$; then $e^{-\beta|u-x|}\le e^{-\beta|x-x'|}e^{\beta|u-x'|}$ and Cauchy--Schwarz give the bound. The quotients follow by replacing one or both ends by their transported differences, $Y$ being linear in each end. No kernel is evaluated at points less than $1$ apart, and the doubly differentiated line enters only through its operator norm.
\end{proof}

\begin{corollary}[the consumed form at $L=5$]\label{sup1:cor:S24}
Let $0\le h,h'\le1$ with $|dh|,|ddh|\le c_h$ (the same for $h'$), the smooth localisers of \cite[(1.33)]{Ba3}; put $d_h\defeq\dist_\infty(\supp h,X_{\mathrm{in}}^c)$, $d_{h'}$ likewise, and assume $d_h,d_{h'}\ge R_0+12$. Then, at any mutual distance of the supports (overlap included),
\[
\norm{h\,\delta\,h'}_{1,\alpha}\le C_{\mathrm{loc}}\exp\!\Big(-\frac{\delta_K}2(d_h+d_{h'}-2R_0-19.4)\Big)e^{-\beta\dist_\infty(\supp h,\supp h')},\qquad C_{\mathrm{loc}}\defeq e^{3\beta}(1+6c_h)^2C_D(\alpha).
\]
Since $d_h+d_{h'}-2R_0-19.4\ge4.6>0$, this is the first alternative of \cite[(2.5)]{Ba3} in the $O(1)$ form, and it carries the depth factor $\exp(-\frac{\delta_K}2(d_h+d_{h'}))$, the form of \cite[(2.26)]{Ba1} with the complement of the smaller set.
\end{corollary}

\begin{proof}
\Cref{sup1:lem:S15} with $\mathcal P\defeq\{\rho\ge d_h-1-\eta\}$, $\mathcal P'\defeq\{\rho\ge d_{h'}-1-\eta\}$, both contained in $\mathcal P_D$ since $R_0+12-1.2=R_0+10.8\ge R_0+10.4$; on $\mathcal P\times\mathcal P'$ \cref{sup1:prop:S23} gives the hypothesis with $M_Y\defeq C_D\exp(-\frac{\delta_K}2[(d_h-1.2-R_0-7)+(d_{h'}-1.2-R_0-10)])$.
\end{proof}

\begin{remark}\label{sup1:rem:S2}
The three instances in \cite{Ba1,Ba2} are $X_{\mathrm{out}}=\Omega$, $X_{\mathrm{in}}=\Omega_2$; $X_{\mathrm{out}}=\Omega$, $X_{\mathrm{in}}=$ the cube; and $X_{\mathrm{out}}=T_\eta$, $X_{\mathrm{in}}=$ the cube, $A=0$. In the first the vertex supports have depth $\ge5r(t)-9$ in $\Omega_2$ (the nesting of \cite[pp.~579--580]{Ba2}: the summations in the vertices are restricted to $\Lambda^{(k)}_7$ while the propagator domain is $B^k(\Lambda^{(k)}_2)$), which exceeds $R_0+12$ under $R\ge3R_0+15$; at $L=5$ the depth $R_0+12$ is uniform in $k\ge1$, there being no separate $k=1$ case ($\eta_1=\frac15$). Near $\partial\Omega_2$ nothing is proved, and nothing follows from (P) alone (Neumann boundary regularity at the edges of a union of large blocks).
\end{remark}
